\documentclass[12pt]{article}
\usepackage{amsmath,amssymb,amsthm}
\usepackage[margin=1in]{geometry}

\newtheorem{theorem}{Theorem}
\newtheorem{lemma}[theorem]{Lemma}
\newtheorem{corollary}[theorem]{Corollary}
\newtheorem{proposition}[theorem]{Proposition}
\theoremstyle{definition}
\newtheorem{definition}[theorem]{Definition}

\title{Project Euler Problem 659: Largest Prime}
\author{Solution}
\date{}

\begin{document}
\maketitle

\section{Problem Statement}

Define $P(k)$ as the largest prime dividing two successive terms of $n^2 + k^2$.
Find the last 18 digits of $\displaystyle\sum_{k=1}^{10^7} P(k)$.

\section{Key Theorem}

\begin{theorem}
$P(k)$ equals the largest prime factor of $4k^2 + 1$.
\end{theorem}

\begin{proof}
If $p \mid \gcd(n^2+k^2, (n+1)^2+k^2)$, then $p \mid (2n+1)$ and $p \mid (n^2+k^2)$.

From $2n+1 \equiv 0 \pmod{p}$: $n \equiv \frac{p-1}{2} \pmod{p}$.

Then $n^2 + k^2 \equiv \frac{(p-1)^2}{4} + k^2 \equiv \frac{1+4k^2}{4} \pmod{p}$.

So $p \mid (4k^2+1)$.

Conversely, for any prime $p \mid (4k^2+1)$, set $n = \frac{p-1}{2}$ to verify $p$ divides two successive terms.

Hence $P(k) = \max\{p : p \text{ prime}, p \mid (4k^2+1)\}$.
\end{proof}

\section{Properties of $4k^2+1$}

\begin{proposition}
Every prime factor $p$ of $4k^2+1$ satisfies $p \equiv 1 \pmod{4}$ or $p = 4k^2+1$ itself is prime.
\end{proposition}

\begin{proof}
$4k^2 \equiv -1 \pmod{p}$ means $-1$ is a quadratic residue mod $p$, so $p \equiv 1 \pmod{4}$ by Euler's criterion (for odd $p > 2$).
\end{proof}

\section{Sieve Algorithm}

\begin{enumerate}
\item Sieve primes $p \equiv 1 \pmod{4}$ up to $\sqrt{4\cdot(10^7)^2+1} \approx 2\times 10^7$.
\item For each prime $p$, find $r$ with $r^2 \equiv -1 \pmod{p}$ (using Tonelli-Shanks).
\item Mark $k$ values: $2k \equiv \pm r \pmod{p}$, i.e., $k \equiv \frac{\pm r}{2} \pmod{p}$.
\item Track $P[k] = \max$ prime factor found.
\item For remaining $k$ where $4k^2+1$ is not fully factored, do primality test.
\end{enumerate}

\section{Answer}

\[
\boxed{238518915714422000}
\]

\end{document}
